%% ============================================================
%%  Chapter 2 – Literature Review
%% ============================================================
\chapter{Literature Review}
\label{ch:lit}

\begin{infobox}[Chapter Overview]
This chapter surveys the state of the art in Medical VQA datasets,
vision-language foundation models, parameter-efficient fine-tuning
strategies, explainability methods, and privacy engineering in medical AI ---
all of which inform the design decisions made in \medxplain.
\end{infobox}

%% ─────────────────────────────────────────────────────────────
\section{Medical Visual Question Answering}
\label{sec:lit_vqa}

Visual Question Answering (VQA) is the task of producing a natural-language
answer $a$ given an image $\mathbf{x}$ and a question $q$~\citep{antol2015vqa}.
The medical subdomain imposes unique challenges: answers are clinically
consequential, question vocabulary is domain-specific, and the mapping from
visual evidence to answer is rarely straightforward.

\subsection{Benchmark Datasets}
\label{subsec:datasets}

Table~\ref{tab:datasets} summarises the principal publicly available
Medical VQA datasets.

\begin{table}[htbp]
\centering
\caption{Publicly available Medical VQA datasets.}
\label{tab:datasets}
\begin{tabular}{@{}lllr@{}}
\toprule
\textbf{Dataset} & \textbf{Domain} & \textbf{Task Type} & \textbf{QA Pairs} \\
\midrule
VQA-Med 2019~\citep{ben2019vqamed}   & Radiology (MedPix)      & Closed/Open & 15{,}292 \\
VQA-Med 2021~\citep{did2021vqamed}   & Radiology               & Closed/Open & 5{,}500  \\
VQA-RAD~\citep{lau2018vqarad}        & Radiology               & Closed/Open & 3{,}515  \\
PathVQA~\citep{he2020pathvqa}        & Histopathology          & Yes/No + Open & 32{,}799 \\
SLAKE~\citep{liu2021slake}           & Radiology (EN + ZH)     & Semantic+Visual & 14{,}028 \\
\midrule
Hospital Cohort (this work)          & Chest X-ray (private)   & Closed      & ---     \\
\bottomrule
\end{tabular}
\end{table}

\noindent
A common weakness of all public benchmarks is their origin in annotated
teaching files rather than prospective clinical workflow. This thesis
addresses this gap by working with real, de-identified hospital data.

\subsection{Evaluation Metrics in Medical VQA}
\label{subsec:metrics_lit}

Standard VQA evaluation uses exact-match accuracy and BLEU~\citep{papineni2002bleu}
scores. Recent work advocates for semantically richer metrics:

\begin{itemize}
  \item \textbf{BERTScore}~\citep{zhang2020bertscore}: Computes token-level
        cosine similarity between candidate and reference using contextual
        BERT embeddings. Better captures semantic equivalence than $n$-gram
        overlap.
  \item \textbf{Clinical Accuracy}: Domain-specific accuracy measuring
        whether the correct clinical entity (finding, diagnosis) is
        identified, independent of surface form.
  \item \textbf{RadGrad Score}~\citep{abacha2019question}: Composite metric
        combining BLEU with clinical relevance.
\end{itemize}

\medxplain{} implements exact-match accuracy, BLEU-4, BERTScore F1, and a
novel Clinical F1 metric based on per-pathology presence detection
(\S\ref{subsec:clinical_f1}).

%% ─────────────────────────────────────────────────────────────
\section{Vision-Language Foundation Models}
\label{sec:lit_vlm}

The field has been transformed by large-scale contrastive pre-training on
paired image-text data, a paradigm pioneered by CLIP~\citep{radford2021clip}.

\subsection{CLIP and Biomedical Variants}
\label{subsec:clip}

CLIP trains dual image and text encoders to maximise cosine similarity for
matching pairs in a batch while minimising it for non-matching pairs,
using the NT-Xent objective:
\begin{equation}
\mathcal{L}_{\mathrm{CLIP}} = -\frac{1}{N} \sum_{i=1}^{N}
\log \frac{\exp(\bm{v}_i \cdot \bm{t}_i / \tau)}
          {\sum_{j=1}^{N} \exp(\bm{v}_i \cdot \bm{t}_j / \tau)}
\label{eq:ntxent}
\end{equation}
where $\bm{v}_i$ and $\bm{t}_i$ are $L_2$-normalised image and text
embeddings for the $i$-th pair, and $\tau$ is a learnable temperature.

\textbf{BiomedCLIP}~\citep{zhang2023biomedclip} extends CLIP by pretraining
on 15 million biomedical image-text pairs harvested from PubMed Central
open-access articles. The resulting Vision Transformer (ViT-B/16) backbone
shows strong transfer to radiology, pathology, and ophthalmology tasks.
\medxplain{} uses BiomedCLIP as its vision encoder.

\subsection{Medical Language Decoders}
\label{subsec:llm}

\textbf{BioGPT}~\citep{luo2022biogpt} is a GPT-2~\citep{radford2019gpt2}
variant pretrained on 15 million PubMed abstracts, achieving state-of-the-art
performance on several biomedical NLP benchmarks including BioASQ and PubMedQA.

More recent work applies instruction-tuned LLMs to radiology
report generation~\citep{lee2023llava,li2023llavamedical}, but these
require substantial compute and typically cannot run in air-gapped hospital
environments. BioGPT balances domain specificity with manageable model size.

\subsection{Multimodal Fusion Architectures}
\label{subsec:fusion}

After CLIP-style alignment, the dominant fusion strategies are:
\begin{enumerate}
  \item \textbf{Concatenation + MLP}: Simple but loses spatial structure.
  \item \textbf{Cross-attention}~\citep{alayrac2022flamingo}: Text tokens
        attend to visual tokens at each transformer layer, enabling rich
        image-conditioned generation.
  \item \textbf{Q-Former}~\citep{li2023blip2}: A lightweight bottleneck
        transformer that extracts task-relevant visual features before
        passing them to the frozen LLM.
\end{enumerate}

\medxplain{} adopts \textbf{cross-attention fusion}, implemented in
\code{medical\_vqa\_infrastructure/models/fusion.py}, because it provides
interpretable attention maps and fits naturally with the Grad-CAM
explainability pipeline.

%% ─────────────────────────────────────────────────────────────
\section{Parameter-Efficient Fine-Tuning}
\label{sec:lit_peft}

Hospital GPU budgets rarely accommodate full fine-tuning of models with
hundreds of millions of parameters. Parameter-Efficient Fine-Tuning (PEFT)
methods update only a small subset of parameters while keeping the
pre-trained backbone frozen.

\subsection{Low-Rank Adaptation (LoRA)}
\label{subsec:lora}

LoRA~\citep{hu2021lora} decomposes the weight update $\Delta W$ of any
linear layer as a product of two low-rank matrices:
\begin{equation}
  W' = W + \Delta W = W + \frac{\alpha}{r}\, B A
  \label{eq:lora}
\end{equation}
where $W \in \mathbb{R}^{d \times k}$ is frozen, $A \in \mathbb{R}^{r \times k}$
and $B \in \mathbb{R}^{d \times r}$ are the trainable adapters (rank $r \ll \min(d,k)$),
and $\alpha$ is a scaling hyperparameter. LoRA initialises $A$ with Kaiming
uniform and $B$ with zeros, so $\Delta W = 0$ at the start of training.

With $r=8$ and applying LoRA to the query, key, and value projections of
the BioGPT attention layers, the number of trainable parameters in the
language decoder drops from ${\sim}347$M to ${\sim}2$M, a reduction of
$99.4\%$.

\subsection{Adapter Layers}
\label{subsec:adapters}

\cite{houlsby2019adapter} introduce learnable bottleneck modules inserted
between sub-layers of each transformer block. While effective, adapters
add inference latency due to sequential computation. LoRA avoids this by
merging $BA$ into the original weight at inference time
($W' = W + \frac{\alpha}{r}BA$), yielding zero additional latency.

%% ─────────────────────────────────────────────────────────────
\section{Explainability in Medical AI}
\label{sec:lit_xai}

Regulatory guidance (EU AI Act, FDA SaMD framework) and clinical common
sense both require that a Medical AI system be able to show \emph{where}
in the image it found evidence for its prediction.

\subsection{Gradient-Weighted Class Activation Maps (Grad-CAM)}
\label{subsec:gradcam_lit}

Grad-CAM~\citep{selvaraju2017gradcam} computes a heatmap by
back-propagating the gradient of the predicted class score $y^c$ with
respect to the feature maps $A^k$ of the final convolutional (or attention)
layer:
\begin{equation}
  \alpha_k^c = \frac{1}{Z} \sum_{i} \sum_{j}
               \frac{\partial y^c}{\partial A^k_{ij}}
  \label{eq:alpha}
\end{equation}
\begin{equation}
  L^c_{\mathrm{Grad\text{-}CAM}} =
  \mathrm{ReLU}\!\left(\sum_k \alpha_k^c A^k\right)
  \label{eq:gradcam}
\end{equation}
The resulting map $L^c$ is resized to the input image resolution and
superimposed as a colour overlay, highlighting regions that most
positively contributed to class $c$.

\medxplain{} computes Grad-CAM for every VQA forward pass, targeting the
last block of the BiomedCLIP ViT, and renders the result in real-time
in the Gradio interface (\S\ref{sec:ui}).

\subsection{Integrated Gradients}
\label{subsec:ig}

Integrated Gradients~\citep{sundararajan2017integratedgradients} provide
a provably axiomatic attribution satisfying completeness and sensitivity,
making it a useful alternative or complement to Grad-CAM:
\begin{equation}
  \mathrm{IG}_i(\mathbf{x}) =
  (x_i - x_i') \int_{\alpha=0}^{1}
  \frac{\partial F(\mathbf{x}' + \alpha(\mathbf{x} - \mathbf{x}'))}
       {\partial x_i}\, d\alpha
  \label{eq:ig}
\end{equation}
where $\mathbf{x}'$ is a baseline (e.g., a black image). Implemented in
\code{explainability/integrated\_gradients.py}.

\subsection{Counterfactual Explanations}
\label{subsec:counterfactual}

\cite{goyal2019counterfactual} demonstrate counterfactual VQA: finding the
minimal change to an image that flips the predicted answer. Implemented in
\code{explainability/counterfactual.py} for offline analysis.

%% ─────────────────────────────────────────────────────────────
\section{Privacy and Compliance in Medical AI}
\label{sec:lit_privacy}

\subsection{HIPAA Safe Harbour De-identification}
\label{subsec:hipaa}

The US Department of Health and Human Services specifies
18 categories of \textsc{PHI}~\citep{hipaa} that must be removed or
generalised for data to be considered de-identified:
names, geographic subdivisions, dates (other than year), phone/fax numbers,
email addresses, social security numbers, medical record numbers,
health plan beneficiary numbers, account numbers, certificate/licence numbers,
vehicle identifiers, device identifiers/serial numbers, URLs, IP addresses,
biometric identifiers, full-face photographs, and any other uniquely
identifying number.

\medxplain{}'s PHI scrubber (\S\ref{sec:phi}) removes all 18 categories
from DICOM headers and employs regex $+$ optional Named Entity Recognition
(NER) for free-text radiology reports.

\subsection{DICOM Confidentiality Profile}
\label{subsec:dicom_profile}

The DICOM Standard Part 15 Annex E (PS~3.15~Annex~E) specifies the
Basic Application Level Confidentiality Profile, enumerating DICOM
attribute tags that must be scrubbed, anonymised, or pseudonymised
before data leaves the institution~\citep{dicom2023}.
\medxplain{} aligns its tag whitelist with this profile (\S\ref{sec:phi}).

\subsection{GDPR and Data Minimisation}
\label{subsec:gdpr}

The General Data Protection Regulation~\citep{gdpr} applies to any
processing of personal data of EU residents. The key principle of
\emph{data minimisation} (Article~5(1)(c)) requires that only the minimum
data necessary for the stated purpose be processed. The \medxplain{}
pipeline enforces this by:
\begin{itemize}
  \item Never logging raw PHI (only SHA-256 file hashes in the audit log).
  \item Setting \code{allow\_cloud\_upload: false} by default.
  \item Stripping source file paths from the clean \code{DICOMStudy} object.
\end{itemize}

\subsection{Federated Learning}
\label{subsec:federated}

Federated Learning~\citep{mcmahan2017federated} trains models across
multiple institutions without exchanging raw data, instead aggregating
gradient updates. While beyond the scope of this thesis, it represents a
natural extension of \medxplain{} for multi-site validation.

%% ─────────────────────────────────────────────────────────────
\section{Summary}
\label{sec:lit_summary}

Table~\ref{tab:related} compares \medxplain{} to representative related
systems along five dimensions.

\begin{table}[htbp]
\centering
\caption{Comparison of \medxplain{} to related Medical AI systems.}
\label{tab:related}
\setlength{\tabcolsep}{5pt}
\begin{tabular}{@{}lccccc@{}}
\toprule
\textbf{System} & \textbf{DICOM} & \textbf{PHI} & \textbf{VQA} & \textbf{XAI} & \textbf{Hospital} \\
 & \textbf{Input} & \textbf{Scrub} & & & \textbf{Deploy} \\
\midrule
CheXNet~\citep{rajpurkar2017chexnet}      & \ding{53} & \ding{53} & \ding{53} & \ding{53} & Partial \\
MAML-VQA~\citep{moon2022mamlvqa}          & \ding{53} & \ding{53} & \ding{51} & \ding{53} & \ding{53} \\
LLaVA-Med~\citep{li2023llavamedical}      & \ding{53} & \ding{53} & \ding{51} & \ding{53} & \ding{53} \\
RGRG~\citep{tanida2023rgrg}               & \ding{53} & \ding{53} & Partial   & Partial   & \ding{53} \\
\midrule
\textbf{\medxplain{} (this work)}         & \ding{51} & \ding{51} & \ding{51} & \ding{51} & \ding{51} \\
\bottomrule
\end{tabular}
\end{table}

\noindent
\ding{51} = fully supported; \ding{53} = not supported; Partial = limited support.
